\section{Conclusion}
  \label{sec:conclusion}

  The O-series has now traversed the full chain from obstruction identification
  to exact measurement.
  O8 identified the geometric obstruction on LPS graphs;
  O9 resolved it via polynomial-growth Heisenberg graphs;
  O10 identified the algorithmic obstruction in the dense TensorSketch;
  O11 lifted this obstruction via the bidimensional Fourier-character proxy,
  obtaining $\hat{\delta}_{\mathrm{cap}} \approx 3.39$;
  O12 closes the observability gap by implementing the exact Weil projection, and
  additionally derives, in closed form, the intrinsic asymptotic exponent that the finite-$q$
  fits of this and the previous papers in the series only ever approximated.

  The central results of this paper are twofold.
  First, the finite-$q$ measurement
  \begin{equation}
    \hat{\delta}_{\mathrm{exact}} \in \{4.37,\,4.28,\,4.80\}
    \qquad (q \in \{29,53,61\},\; R^2 > 0.98),
    \label{eq:conclusion-finite-q}
  \end{equation}
  independently reproduced by pure combinatorial set arithmetic (Section~\ref{sec:results}).
  Second, and decisively, Section~\ref{sec:comparison} shows these values are finite-$q$
  crossover slopes of the block saturation-depth distribution, not estimates of an intrinsic
  asymptotic exponent, and derives that exponent exactly:
  \begin{equation}
    \delta_{\mathrm{exact}} = 3
    \qquad \bigl(q\to\infty,\ n\to\infty,\ n=o(\sqrt{q})\bigr),
    \label{eq:conclusion-asymptotic}
  \end{equation}
  from Heisenberg shell growth alone, holding block-by-block (not only in the mean) with
  probability $1-O(n/q)$.
  This exponent is exactly independent of the central coordinate $\gamma$ of
  $\mathrm{Heis}_3(\mathbb{Z}/q\mathbb{Z})$ and of the metaplectic phase $\psi_c(\gamma)$:
  Proposition~\ref{prop:alpha-beta} and Remark~\ref{rem:blind} show the rank observable
  $\Sigma_n^{(c)}$ depends on the visited group elements only through their $b$-coordinates,
  so no coherence mechanism built from $\gamma$ can affect it, by algebraic identity rather
  than empirical regularity.

  The result carries both a positive and a cautionary message.
  On the positive side, a measurable and now exactly characterised power-law decay of
  $\Sigma_n^{(c)}$ is established, block-by-block, in the unfolded regime, with $R^2 > 0.98$
  confirmed at every finite-$q$ prime tested and the asymptotic exponent proved rather than
  fitted.
  On the cautionary side, the accessible finite-$q$ range ($q \leq 61$) sits inside the
  crossover regime rather than the unfolded one: the fitted slopes there are real,
  reproducible statistics of the same construction, but they are not, and were not intended
  to be read as, the asymptotic exponent itself.

  No constraint on the cascade exponent $\beta^*$ is drawn from the present paper.
  Section~\ref{sec:beta} shows that substituting either the finite-$q$ crossover slopes or
  the unfolded exponent $\delta_{\mathrm{exact}}=3$ into the O3--O7 relation
  $\beta^*=1/(\delta+\tfrac12)$~\cite{Beau2026a11} would presuppose an answer to a question
  this paper does not settle: which of the regimes identified in
  Section~\ref{sec:three-regimes} that relation's derivation takes as input.
  Resolving this requires re-auditing O3--O7 itself.

  O12 is the first point in the O-series where the Cosmochrony framework is confronted with
  exact representation-level data, and the first where the asymptotic capacity exponent of
  this construction is derived rather than fitted.
  The resulting picture sharpens the programme: the central coordinate $\gamma$, contrary to
  the motivation for building the exact projection, is provably inert for this observable,
  and the finite-$q$ numbers previously reported carry no evidence about it either way.
  The next steps are to re-audit which regime the O3--O7 structural relation legitimately
  takes as input (Section~\ref{sec:beta}), and to extend the classification of
  Section~\ref{sec:comparison} to the O13 and O15 constructions that share this pipeline.
